% Point-by-Point Response to Reviewers
% Manuscript: "Universal spectral scaling of terrain-induced fatigue"
% Journal: Nature Communications
% Handling Editor: Dr. Shuyan Xiong
%
% Conventions used in this document:
%   - Reviewer comments are quoted verbatim in italics
%   - Our responses follow immediately, in plain text
%   - "Revised manuscript, [location]" directs the reviewer to the specific change
%   - New or changed text from the manuscript is shown inline in \textbf{bold}

\documentclass[11pt]{article}
\usepackage[margin=1in]{geometry}
\usepackage{amsmath,amssymb}
\usepackage{xcolor}
\usepackage{hyperref}
\usepackage{setspace}
\usepackage{enumitem}
\onehalfspacing

\newcommand{\revquote}[1]{\medskip\noindent\textit{``#1''}\medskip}
\newcommand{\revloc}[1]{\noindent\textbf{Location in revised manuscript:} #1\medskip}

\begin{document}

\begin{center}
{\large\bfseries Point-by-Point Response to Reviewers}\\[6pt]
\textit{Spectral scaling of natural-terrain-induced fatigue}\\[4pt]
Nature Communications\\[4pt]
\end{center}

\noindent We thank the three reviewers and the handling editor for their careful and constructive reading of the manuscript. The comments collectively identify real gaps that we have addressed through manuscript revision, additional analysis, and clarification of scope. We respond to each comment in turn below, quoting the reviewer verbatim and providing our response. All changes in the revised manuscript are highlighted in blue.

\bigskip\noindent\rule{\textwidth}{0.4pt}

% ==============================================================
\section*{Reviewer \#1}
% ==============================================================

\noindent Reviewer \#1 raised six technical issues and three formatting/presentation issues. We address each below.

\bigskip

% --------------------------------------------------------------
\subsection*{Technical Issue (1): ISO 8608 comparison}
% --------------------------------------------------------------

\revquote{The manuscript establishes a fractal-based framework. However, a straightforward approach would be to directly generate road profiles from ISO 8608 and then use them as input to excite a quarter-car model for calculating vibration energy and fatigue damage. The authors should clarify what quantifiable advantages the proposed fractal framework offers over the direct use of ISO 8608 parameters, in terms of prediction accuracy, scope of applicability, or practical operability.}

\noindent\textbf{Response.} We thank the reviewer for raising this point. In retrospect, we did not adequately distinguish the proposed framework from the ISO~8608 direct-simulation pipeline, and the reviewer is correct that this needed to be addressed explicitly. We have revised the manuscript in two places.

The structural parallel between the two approaches is real and should be acknowledged: ISO~8608 models road PSD as $G_d(n) = G_d(n_0)(n/n_0)^{-w}$, where $G_d(n_0)$ is a spectral amplitude and $w$ is a waviness exponent. Our $(C_z,\beta)$ framework uses identical two-parameter PSD logic, with $C_z \leftrightarrow G_d(n_0)$ and $\beta \leftrightarrow w$. The difference lies not in the mathematical form but in what each approach can do \emph{before} a vehicle has traversed the terrain.

We identify three concrete advantages, the first two of which are straightforwardly quantifiable; the third describes improved explanatory power relative to amplitude-only roughness descriptors rather than a direct head-to-head ISO~8608 benchmark:

\begin{enumerate}[leftmargin=*, label=(\roman*)]

\item \textbf{Predictive capability from remote sensing alone.}
ISO~8608 requires a profilometer measurement to determine $w$ for a given surface; in the absence of measurement data, the standard fixes $w \approx 2$ as an empirical constant. Our framework derives $\beta$ from terrain geometry via $\beta = 7-2D$, so spectral slope can be predicted from any digital elevation model without in-situ instrumentation. This is the primary practical advantage for the intended application domain: off-road military operations, planetary rover missions, and pre-deployment route planning where profilometer access is unavailable.

\item \textbf{Scope: natural terrain beyond engineered road classes.}
ISO~8608 road classes (A--H) were calibrated on paved and graded surfaces. They do not extend to natural terrain (mountains, desert washes, off-road military routes) where no road classification exists. Our framework was specifically developed and validated for this domain (13 USGS 3DEP LiDAR regions spanning flat plains to high-relief mountain terrain, $r = -0.62$, $p < 0.05$). We explicitly note that the framework is \emph{not} intended to replace ISO~8608 for maintained paved roads where profilometer data are available; it extends predictive capability to the domain ISO~8608 does not cover.

\item \textbf{Improved explanatory power relative to amplitude-only roughness descriptors.}
Standard ISO~8608 practice classifies roads by amplitude alone, fixing $w \approx 2$. Explicitly fitting both parameters ($C_z$ and $\beta$) improves prediction of vibration energy from $R^2 = 0.91$ (amplitude-only) to $R^2 = 0.96$ in our simulations. This improvement represents a measurable reduction in residual variance that propagates through the nonlinear S--N relationship and can materially affect fatigue-life estimates, given that $N_f \propto E^{-m/2}$ with $m$ typically ranging from 3 to 10 for relevant materials.

\end{enumerate}

\noindent A new comparison table (Table~2 in the revised manuscript) summarises these distinctions side by side. We note that the $R^2 = 0.91$ entry in the table refers to the amplitude-only single-parameter baseline from our own simulations (Fig.~4d), not to a dedicated ISO~8608 simulation; the table label makes this explicit (``amplitude-only baseline'') to avoid any implication of a direct head-to-head comparison with ISO~8608 profilometer data.

We have also moderated language around fatigue claims throughout the manuscript. Where the original text claimed ``quantitative fatigue predictions'' or ``accurate fatigue prediction'' from DEM data alone, the revised text uses ``quantitative vibration-severity estimates,'' ``pre-deployment fatigue-risk proxies,'' or ``accurate prediction of vibration severity and fatigue risk'': language consistent with the actual chain of evidence (terrain PSD $\rightarrow$ vibration energy $\rightarrow$ fatigue-damage proxy), without implying a validated absolute fatigue-life number. The phrase ``enabling accurate fatigue prediction for any vehicle given its natural frequency'' (extended vehicle ensemble section) has similarly been revised to ``enabling estimation of relative fatigue-risk trends across vehicle classes with known natural frequency.'' The $R^2 = 0.91 \rightarrow 0.96$ improvement is described as a ``measurable'' rather than ``substantial'' reduction in residual variance, consistent with the reviewer's own observation that the 5-point gain may not be a game-changing improvement in isolation.

We have also taken care not to overstate the comparison. The $R^2$ improvement from 0.91 to 0.96 comes from existing simulation data (Fig.\ 4d in the original manuscript) comparing single-parameter and two-parameter fits; it is not a new dedicated simulation. Furthermore, the DEM-based predictions are validated on natural terrain but explicitly shown to fail on flat paved urban roads (Copenhagen, Fig.\ 6), consistent with the domain boundaries described above.

\revloc{Main manuscript: Introduction, paragraph 4 (new ISO~8608 paragraph, highlighted in blue); Discussion, new subsection ``Relationship to ISO~8608 and prediction accuracy'' with Table~2. Supplementary: no changes for this comment.}

\bigskip

% --------------------------------------------------------------
\subsection*{Technical Issue (2): Vehicle independence vs.\ frequency dependence (Figures 3 and 5)}
% --------------------------------------------------------------

\revquote{Figure 3 uses only three vehicle examples (A, B, and C) with similar frequencies, and based on this, the manuscript concludes that the results are only weakly influenced by vehicle selection. However, the results shown in Figure 5 appear to be frequency-dependent. Taken together, Figures 3 and 5 do not support the claim that the proposed framework is largely insensitive to vehicle selection. The statement in the abstract that the framework is vehicle-independent may therefore be misleading.}

\noindent\textbf{Response.} We agree with the reviewer. The original manuscript overstated vehicle independence, and the contradiction between Figures~3 and~5 is real. We have revised the manuscript throughout to remove the claim of vehicle independence and replace it with the more accurate characterisation: \emph{terrain-dominated, with predictable frequency-dependent modulation}.

The distinction the reviewer identifies is precise: Figure~3 uses three vehicles spanning only a narrow natural frequency range (0.89--0.98 Hz). Within that narrow band, vehicle type explains only 1.9\% of energy variance, and the scaling exponent varies by less than 2\% (CV = 1.4\%). This is accurately described as \emph{weak} vehicle sensitivity within the tested range, not as vehicle independence in general. Figure~5, by contrast, spans 100 vehicles from 0.80 to 2.49 Hz, and shows clear systematic frequency dependence ($\gamma = 0.36f_n - 2.93$, $R^2 = 0.879$). These two results are fully consistent with each other: the apparent vehicle independence in Figure~3 is a consequence of the narrow frequency range tested, not a general property of the framework. This is now made explicit in the manuscript.

Specific changes made:

\begin{itemize}[leftmargin=*]
\item Abstract: ``vehicle configuration explains only 2.0\%'' is now qualified as ``within the three-vehicle narrow-frequency comparison (0.89--0.98 Hz); broader ensemble simulations reveal systematic natural-frequency modulation.'' The updated value from the data-verified regression is 1.9\%.
\item Introduction: ``vehicle-independent scaling law'' replaced with ``terrain-dominated scaling law \ldots with natural frequency providing predictable modulation.''
\item Results (energy scaling section): ``vehicle-independent scaling'' replaced with ``terrain-dominated scaling \ldots within the narrow frequency range 0.89--0.98 Hz.''
\item Results (variance decomposition): ``vehicle-independent terrain severity index'' replaced with ``terrain-dominated severity index \ldots with the caveat that natural frequency modulates the scaling exponent predictably across a broader vehicle ensemble.''
\item Discussion opening: ``vehicle-independent scaling'' replaced with ``terrain-dominated scaling.''
\item Discussion: ``largely vehicle-independent predictions'' replaced with ``terrain-dominated predictions with predictable frequency-dependent modulation once vehicle natural frequency is known.''
\item Universality paragraph: removed claim that the exponent is ``remarkably insensitive to vehicle parameters''; replaced with explicit acknowledgement that the narrow-frequency result and broad-ensemble result are complementary.
\end{itemize}

No new simulations are required to address this comment. The evidence needed to support the revised framing is already present in Figures~3 and~5; the issue was one of interpretation and language, not missing data.

\revloc{Main manuscript: Abstract (sentence 3, variance numbers and frequency qualification); Introduction paragraph 6 (``terrain-dominated scaling law''); Results ``Energy scaling'' final paragraph; Results ``Terrain-dominated scaling and multi-vehicle validation'' (Table~1 values, variance decomposition text, ``These results'' paragraph). Discussion opening sentence. All highlighted in blue. Supplementary: Note 4 (MI conclusion), Note 5 (profile bias), Note 6 (framework implications), Note 10 (section title, intro, results text, implications subsection). All highlighted in blue.}

\bigskip

% --------------------------------------------------------------
\subsection*{Technical Issue (3): Slope discrepancy between Figures 2 and 3}
% --------------------------------------------------------------

\revquote{The fitted slopes in Figures 2 and 3 differ by a factor of two ($-1.59$ vs.\ $-3.18$), yet the two fittings appear to correspond to the same regression, since $R^2$, $r$, and $n$ are exactly identical. This inconsistency needs to be explained.}

\noindent\textbf{Response.} The reviewer correctly identifies a genuine inconsistency in the original manuscript. We thank them for catching this. The two slopes describe \emph{different regressions}, but the original text and Figure~3 legend did not make this clear; in fact, Figure~3 panel (a) was mislabelled.

The correct interpretation is as follows:

\begin{itemize}[leftmargin=*]
\item \textbf{Figure~2, slope $-1.59$:} regression of the vehicle acceleration \emph{spectral exponent} $\beta_a$ on fractal dimension $D$, i.e., $\beta_a = -1.59 \cdot D + 4.69$ ($R^2 = 0.913$, $r = -0.956$, $n = 1500$).

\item \textbf{Figure~3 panel (a), slope $-3.11$:} regression of log \emph{vibration energy} $\log E$ on $D$, i.e., $E \propto (D-2)^{-3.11}$ ($R^2 = 0.915$, $n = 1500$). This is the energy--fractal-dimension scaling, consistent with the amplitude-complexity coupling $E \propto (D-2)^{-3.1}$ reported throughout the paper.
\end{itemize}

The coincidence of identical $R^2$, $r$, and $n$ values across both figures in the original text was a transcription error: the Figure~3 legend mistakenly carried over the statistics from Figure~2 rather than those from the actual energy regression. The underlying physics is consistent: $\beta_a$ vs $D$ has slope $-1.59$ (spectral slope, suspension-filtered, $R^2 = 0.913$); $\log E$ vs $D$ has slope $-3.11$ (energy scaling, $R^2 = 0.915$). These are different regressions with slightly different goodness-of-fit. In the revised Figure~3 legend we report only the energy regression statistic ($R^2 = 0.915$) and explicitly distinguish it from Figure~2, removing the copied $r$ and confidence interval values that belonged to the spectral-exponent regression.

We have corrected the Results text, Figure~3 caption, and Figure~3 legend to clearly identify panel (a) as the \emph{energy--fractal-dimension} plot, distinguish it explicitly from the spectral-slope regression in Figure~2, and explain why the two slopes differ.

\revloc{Main manuscript: Results ``Energy scaling law'' paragraph beginning ``Figure~3 validates...''; Figure Legends, Figure~3 panel~(a) description. Both highlighted in blue.}

\bigskip

% --------------------------------------------------------------
\subsection*{Technical Issue (4): Overgeneralization in title and abstract}
% --------------------------------------------------------------

\revquote{The manuscript states that the proposed framework is not applicable to paved roads. If so, the wording in the title and abstract should be tightened accordingly to avoid overgeneralization.}

\noindent\textbf{Response.} We agree. The manuscript explicitly demonstrates that the DEM-based methodology fails on flat paved urban roads (Copenhagen, Figure~6), making the original title ``Universal spectral scaling of terrain-induced fatigue'' an overgeneralization. We have made two changes:

\begin{enumerate}[leftmargin=*]
\item Title revised from ``Universal spectral scaling of terrain-induced fatigue'' to ``Spectral scaling of natural-terrain-induced fatigue.'' This directly scopes the claim to the domain where the framework is validated.
\item Abstract final sentence revised to explicitly state the domain boundary: the framework applies to ``diverse natural terrain conditions where digital elevation models are available'' and ``is not intended for engineered paved surfaces where pavement micro-texture dominates vehicle excitation.''
\end{enumerate}

\noindent No new simulations are needed; this is a framing correction.

\revloc{Main manuscript: Title; Abstract, final sentence. Supplementary: title page. All highlighted in blue.}

\bigskip

% --------------------------------------------------------------
\subsection*{Technical Issue (5): Whether $C_z$ alone is sufficient (exponent on $\beta_a$ is only $-0.09$)}
% --------------------------------------------------------------

\revquote{The manuscript claims that the two-parameter framework improves upon the one-parameter framework, but the reported improvement is only about 5\%. Moreover, the exponent of $\beta_a$ is only $-0.09$, which is numerically close to zero. This raises the question of whether $C_z$ alone may already serve as a sufficient predictor.}

\noindent\textbf{Response.} This is a fair observation. In the original coupled dataset, the near-zero exponent on $\beta_a$ ($-0.09$) and partial $R^2 = 0.005$ after controlling for $C_z$ (Supplementary Figure~3) do suggest that $C_z$ alone captures most energy variance. However, this is an artifact of strong collinearity ($r = -0.962$) between $C_z$ and $\beta$ in the Diamond-Square terrain generation, not evidence that spectral slope is physically unimportant.

To directly test this, we conducted a \textbf{new decoupling experiment}: 450 simulations (5 fractal dimensions $\times$ 30 realizations $\times$ 3 vehicles) with terrain profile RMS amplitude held exactly constant (CV $< 10^{-14}$\%) while varying only spectral structure through $D$. To be explicit: in the original coupled dataset, $C_z$ is the dominant empirical predictor and $\beta$ has near-zero residual effect after controlling for $C_z$ (partial $R^2 = 0.005$). The constant-amplitude experiment was therefore added specifically to test whether spectral structure has an independent physical effect when amplitude is controlled. The result is definitive:

\begin{itemize}[leftmargin=*]
\item Energy varies strongly with $D$ at constant amplitude: $\log_{10} E = 2.96 \cdot D - 5.49$, $R^2 = 0.834$, $p < 10^{-176}$, $n = 449$.
\item Energy \emph{increases} with $D$ (positive slope = +2.96), opposite to the coupled case (slope = $-3.11$), because higher $D$ places more spectral energy near the suspension resonance.
\item Results are consistent across all three vehicles ($R^2 = 0.82$--$0.85$).
\item $\beta$ vs $D$ holds at constant amplitude (slope = $-1.82$, $R^2 = 0.79$), confirming spectral structure varies as predicted by $\beta = 7-2D$.
\end{itemize}

\noindent This demonstrates that the near-zero $\beta$ exponent in the original regression is a collinearity artifact, not a physical irrelevance. When amplitude is controlled, terrain spectral structure (mediated by fractal dimension through $\beta = 7-2D$) independently explains 83\% of energy variance through the spectral filtering mechanism. The two-parameter framework is justified: $C_z$ controls amplitude, $\beta$ controls frequency content relative to the vehicle resonance, and both are needed for complete terrain characterization.

\revloc{Main manuscript: Discussion, new paragraph ``Role of spectral slope $\beta$ in the two-parameter model'' (substantially expanded with new decoupling results, highlighted in blue). Supplementary: new Note~12 ``Constant-Amplitude Decoupling Experiment'' with full methods, results, and Supplementary Figure~5. Supplementary Figure~3 caption revised (blue).}

\bigskip

% --------------------------------------------------------------
\subsection*{Technical Issue (6): Single speed (15 m/s)}
% --------------------------------------------------------------

\revquote{The simulations are carried out at only one vehicle speed, 15 m/s. However, existing vehicle dynamic studies indicate that different speeds may couple with road conditions to different extents. This issue should be discussed.}

\noindent\textbf{Response.} We agree that speed dependence is an important consideration and have added explicit discussion of this limitation. The analytical framework already predicts how speed enters: from Equation~5, $E \propto v^{6-2D}$, indicating that (i)~energy increases with speed, and (ii)~the speed sensitivity itself depends on fractal dimension: smoother terrain (lower $D$) shows stronger speed dependence. This coupling means that absolute energy values apply strictly at 15~m/s, though the spectral scaling relationships ($\beta = 7-2D$) are speed-independent because they describe PSD shape rather than magnitude. We note that at substantially different speeds, relative terrain severity rankings could shift if fractal dimensions differ, and we identify multi-speed validation (5--30~m/s) as a priority for future work.

No new simulations are required to address this comment because the analytical speed dependence is already derived in the framework (Equation~5). However, experimental multi-speed validation would strengthen the work and is now explicitly listed in Future Work.

\revloc{Main manuscript: Discussion, new paragraph ``Speed dependence'' (highlighted in blue); Future Work, first sentence identifies multi-speed simulation as priority (blue). Supplementary: no changes for this comment.}

\bigskip

% ==============================================================
\subsection*{Formatting and Presentation Issues}
% ==============================================================

\revquote{(1) There are erroneous citation symbols ([?]) at Lines 331 and 731.}

\noindent\textbf{Response.} These arose from missing bibliography entries for two citations (Cover \& Thomas 2006; Musgrave et al.\ 1989) that were referenced in the text but not included in the bibliography. Both are now added to the reference list and resolve correctly on compilation.

\revloc{Bibliography: added \texttt{Cover2006} and \texttt{Musgrave2000} entries.}

\bigskip

\revquote{(2) Figure 1 does not contain the labels (a), (b), (c), and (d).}

\noindent\textbf{Response.} Panel labels (a)--(d) have been added to the Figure 1 image file to match the caption references.

\revloc{Figure file: \texttt{figures/Figure1.png} updated with panel labels.}

\bigskip

\revquote{(3) In the Introduction, the statement ``Military vehicle maintenance accounts for 60--70\% of total lifecycle costs'' requires an appropriate reference.}

\noindent\textbf{Response.} We have softened the claim to remove the specific unsourced percentages. The revised text reads: ``Military vehicle maintenance can account for a large fraction of total lifecycle costs'' and cites Wong (2022) as a general reference for military ground vehicle maintenance considerations. The specific ``3--5 times shorter'' trucking claim has similarly been softened to ``substantially shorter.''

\revloc{Main manuscript: Introduction, paragraph 2 (highlighted in blue).}

\bigskip

% ==============================================================
\section*{Reviewer \#2}
% ==============================================================

\noindent Reviewer \#2 raised four major comments and numerous minor comments. We address the major comments below.

\bigskip

% --------------------------------------------------------------
\subsection*{Major Comment 1 --- Overstated scope (``universal,'' paved roads, fatigue validation)}
% --------------------------------------------------------------

\revquote{The manuscript overstates its scope. The term ``universal'' is not supported, as the framework is explicitly shown to fail for paved roads where excitation is governed by micro-texture. The applicability is therefore restricted to rough/natural terrain, contradicting the title and parts of the abstract and discussion. In addition, the use of ``terrain-induced fatigue'' suggests a validated fatigue-life prediction framework, which is not demonstrated.}

\noindent\textbf{Response.} We agree on both points and have made substantial revisions accordingly.

\textbf{Scope restriction.} The title has been changed from ``Universal spectral scaling of terrain-induced fatigue'' to ``Spectral scaling of natural-terrain-induced fatigue,'' directly scoping the claim to the domain where the framework is validated. The abstract now ends with an explicit domain statement: the framework applies to ``natural terrain where digital elevation models are available'' and ``is not applicable to engineered paved surfaces.'' The Discussion retains the Copenhagen negative result (Figure~6) as a scientifically valuable boundary condition rather than hiding it.

\textbf{Fatigue-life overclaim.} We acknowledge that the manuscript does not demonstrate a calibrated, absolute fatigue-life prediction. The stress transfer parameter $K_\sigma$ is component-specific and uncalibrated in this work; the S--N curve parameters are generic structural steel values. What the framework demonstrates is the \emph{scaling relationship} between terrain spectral properties and vibration energy, and that this scaling propagates through classical fatigue mechanics with predictable exponents. Accordingly, we have revised all claims from ``fatigue prediction'' to ``vibration-severity estimation and fatigue-risk scaling'':

\begin{itemize}[leftmargin=*]
\item ``quantitative fatigue predictions'' $\rightarrow$ ``quantitative vibration-severity estimates and pre-deployment fatigue-risk proxies''
\item ``accurate fatigue prediction'' $\rightarrow$ ``accurate prediction of vibration severity and fatigue risk''
\item ``enables predictive maintenance'' $\rightarrow$ framed as enabled by the scaling relationships, not by absolute life numbers
\end{itemize}

\noindent The title retains ``fatigue'' because the paper does demonstrate fatigue \emph{scaling} (the 2:1 exponent ratio validating Basquin's law, the 204-fold damage range), even though it does not demonstrate calibrated fatigue-life prediction for a specific component. We believe this is an appropriate and defensible use of the term.

No new simulations are required. The changes are in framing and language.

\revloc{Title; Abstract final sentence; Discussion (multiple paragraphs); Results text where fatigue claims appeared. All instances changed consistently.}

\bigskip

% --------------------------------------------------------------
\subsection*{Major Comment 2 --- Fatigue workflow is simplified and uncalibrated}
% --------------------------------------------------------------

\revquote{Although a fatigue workflow is introduced (rainflow counting, S--N curves, Miner's rule), it is based on highly simplified and generic assumptions, with no explicit component, material, or load-cycle definition. Mean stress effects are neglected, and no load ratio is specified. Furthermore, the stress transfer parameter $K_\sigma$ is not calibrated and is acknowledged to require component-specific identification. As a result, the model does not provide a physically grounded fatigue life prediction and claims regarding fatigue life estimation or maintenance scheduling are not sufficiently supported.}

\noindent\textbf{Response.} The reviewer is correct, and we accept this criticism. The fatigue workflow in this manuscript is intentionally simplified to serve a specific purpose: demonstrating that terrain-driven vibration scaling propagates through classical fatigue mechanics with predictable exponents (the 2:1 ratio between fatigue and energy exponents, validating Basquin's law). It does not constitute a calibrated fatigue-life prediction for any specific component.

We have added an explicit ``Scope of fatigue analysis'' paragraph in the Methods section that states this unambiguously:

\begin{itemize}[leftmargin=*]
\item $K_\sigma$ is uncalibrated and component-specific
\item The S--N exponent $m = 5$ is generic structural steel, not material-specific
\item Mean stress effects and load ratio ($R$) are neglected
\item Miner's rule is a first-order linear accumulation model
\item These simplifications are appropriate for establishing \emph{scaling relationships} but not for absolute life prediction
\end{itemize}

\noindent We have also revised all claims throughout the manuscript from ``fatigue prediction'' and ``maintenance scheduling'' to ``fatigue-risk scaling'' and ``vibration-severity estimation'' (as detailed in our response to Major Comment~1).

The scientific contribution regarding fatigue is the demonstration that:
\begin{enumerate}[leftmargin=*]
\item The fatigue scaling exponent $\gamma_f$ is precisely twice the energy exponent $\gamma_E$ (ratio = $2.00 \pm 0.01$ across 100 vehicles), validating that $\sigma \propto \sqrt{E}$ and Basquin's law ($m = 4$) are internally consistent.
\item The terrain-to-fatigue causal chain preserves scaling relationships at each step.
\item Relative fatigue damage varies by 204-fold across the fractal dimension range, demonstrating that terrain spectral properties have operationally significant consequences even under simplified assumptions.
\end{enumerate}

\noindent These are scaling results, not absolute life predictions. We believe this distinction is now clear in the revised manuscript. Calibrated component-level fatigue prediction is identified as future work requiring FEA-based $K_\sigma$ determination and material-specific S--N testing.

No new simulations are required.

\revloc{Methods, new paragraph ``Scope of fatigue analysis''; throughout manuscript where fatigue claims were softened (same locations as Major Comment~1 response).}

\bigskip

% --------------------------------------------------------------
\subsection*{Major Comment 3 --- Fatigue not consistently integrated in Results}
% --------------------------------------------------------------

\revquote{Fatigue is presented as a central outcome in the title and abstract, but it is not consistently developed in the Results section. Instead, the fatigue formulation appears mainly in the Methods section and is not convincingly linked to the preceding analysis. This disconnect suggests that fatigue is not fully integrated into the core contribution, but rather appended as a post hoc extension, weakening the overall coherence of the manuscript.}

\noindent\textbf{Response.} We understand this perception and have improved the structural coherence. In the original manuscript, the logical flow from vibration energy to fatigue was not signposted clearly enough. In the revised version:

\begin{enumerate}[leftmargin=*]
\item The Results subsection ``Complete fatigue life scaling law'' now opens with an explicit bridging paragraph that states: the preceding sections established terrain $\rightarrow$ vibration energy; this section completes the chain to fatigue damage.

\item The ``Extended vehicle ensemble'' section (Figure~5) reports fatigue results directly in Results, not Methods: the 2:1 ratio between $\gamma_f$ and $\gamma_E$, the mass/frequency independence/dependence of the fatigue exponent, and the 204-fold damage range. These are simulation results, not methodology.

\item The Methods section contains only the \emph{computational procedure} (rainflow counting algorithm, S--N curve parameters, Miner's rule implementation). A new ``Scope of fatigue analysis'' paragraph in Methods now explicitly distinguishes what is methodology (the fatigue computation procedure) from what is a result (the scaling relationships that emerge from it).
\end{enumerate}

\noindent We acknowledge that the paper's core strength is the terrain-to-vibration chain ($D \rightarrow \beta \rightarrow E$), and that the fatigue extension is less rigorously validated than the spectral analysis. However, we believe fatigue belongs in the Results because: (a)~the 2:1 exponent ratio is an empirical finding from the simulations, not an assumption; (b)~the 204-fold damage range quantifies operational significance; and (c)~the fatigue scaling exponent's frequency dependence ($\gamma = 0.36f_n - 2.93$) is a new physical result. The title uses ``fatigue'' to refer to this scaling chain, not to claim calibrated life prediction---a distinction now made explicit throughout the manuscript.

No new simulations are required. The changes are structural and presentational.

\revloc{Results subsection ``Complete fatigue life scaling law,'' new bridging paragraph; Methods, new ``Scope of fatigue analysis'' paragraph; Supplementary Note~3, new scope statement.}

\bigskip

% --------------------------------------------------------------
\subsection*{Major Comment 4 --- Validation insufficient}
% --------------------------------------------------------------

\revquote{The validation is not sufficient to support the claimed impact: the study relies heavily on synthetic terrain with known artifacts, while validation using real data is limited in size and only partially consistent with the theoretical framework.}

\noindent\textbf{Response.} We accept this criticism as the most substantive limitation of the manuscript. The validation strategy has been strengthened through additional analysis and more honest framing of what the current evidence supports.

\textbf{Summary of current validation evidence:}

\begin{enumerate}[leftmargin=*]
\item \textbf{Synthetic terrain (1500 + 18,000 simulations):} Validates the theoretical framework $\beta_t = 7-2D$ and the scaling cascade terrain $\rightarrow$ vibration $\rightarrow$ fatigue under controlled conditions. The Diamond-Square amplitude-complexity coupling is now explicitly acknowledged as an artifact (not a physical law), and the constant-amplitude decoupling experiment ($R^2 = 0.834$, $n=449$) demonstrates that spectral structure independently affects vibration energy.

\item \textbf{USGS 3DEP LiDAR (13 terrain tiles):} Real terrain validation of the $D \rightarrow \beta_t$ relationship yields weighted $r = -0.62$ ($p < 0.05$), confirming that the negative relationship predicted by theory exists in independently measured natural terrain. The correlation is moderate (not strong), consistent with real terrain departing from ideal fBm due to multiple geological processes. Expansion to 25 class-balanced regions is in progress.

\item \textbf{LiRA-CD vehicle sensors (8,609 road segments):} Real vehicle validation of the $\beta_a \rightarrow E$ relationship on a large sample. Both $\beta_a$ vs IRI ($r = +0.194$, $p = 9.8 \times 10^{-74}$) and $E$ vs IRI ($r = +0.145$, $p = 8.6 \times 10^{-42}$) are highly significant, confirming that spectral exponent and vibration energy correlate with independently measured road roughness on 8,609 real road segments.

\item \textbf{Copenhagen negative result:} Demonstrates the methodology's domain boundary---DEM-based characterization does not apply to paved roads where micro-texture dominates. This is a scientifically valuable boundary condition, not a failure.
\end{enumerate}

\textbf{What the validation does and does not support:}

We acknowledge that the validation supports the \emph{existence and direction} of the predicted relationships on real data, but does not yet demonstrate quantitative prediction accuracy sufficient for operational deployment. Specifically:

\begin{itemize}[leftmargin=*]
\item The $D \rightarrow \beta_t$ relationship is confirmed in sign and significance on real terrain, but with moderate correlation ($r = -0.62$) rather than the near-perfect correlation ($r = -0.956$) observed in synthetic terrain. This gap is expected (real terrain $\neq$ ideal fBm) but means the framework provides a first-order approximation, not a precision tool, for natural terrain characterization.

\item The $\beta_a \rightarrow E$ relationship is confirmed with very high statistical power ($n = 8{,}609$, $p < 10^{-40}$), but the effect sizes are modest ($r \approx 0.15$--$0.19$), reflecting the many confounding factors in real driving (speed variation, non-terrain vibration sources, GPS matching uncertainty).

\item The complete end-to-end chain (DEM $\rightarrow$ $D$ $\rightarrow$ $\beta_t$ $\rightarrow$ vehicle $\beta_a$ $\rightarrow$ $E$ $\rightarrow$ fatigue) has not been validated on a single real dataset with all links measured simultaneously.
\end{itemize}

\textbf{Revised framing:} We have revised the manuscript to clearly distinguish between what is demonstrated (scaling relationships, direction of effects, statistical significance) and what remains to be validated (quantitative prediction accuracy for operational use). Claims have been moderated throughout from ``enables accurate fatigue prediction'' to ``enables quantitative vibration-severity estimates and fatigue-risk proxies.'' The abstract and Discussion now explicitly frame the real-data validation as preliminary evidence supporting the theoretical framework, not as operational proof.

\textbf{Additional validation in progress:} We are preparing analysis of the TartanDrive off-road autonomous driving dataset (Carnegie Mellon, 2022), which contains synchronized IMU accelerometer data (200~Hz) and LiDAR-derived terrain heightmaps from an all-terrain vehicle operating at speeds comparable to our simulations ($\sim$15~m/s) on natural off-road terrain. This dataset would enable the first simultaneous measurement of terrain fractal dimension and vehicle vibration response on real off-road terrain. Results will be included in the final revision if data processing is completed in time; otherwise, this is identified as immediate future work.

\textbf{UPDATE --- TartanDrive off-road validation completed.} We have now analyzed the TartanDrive dataset (Triest et al., ICRA 2022): a Yamaha Viking ATV traversing diverse natural terrain (forest trails, gravel paths, open fields) with IMU at 200~Hz and GPS/odometry at 50~Hz. Across 42 trajectory segments (9 traversals, 10-second windows):

\begin{itemize}[leftmargin=*]
\item Vibration energy correlates with spectral slope: $r = 0.512$, $p = 5.4 \times 10^{-4}$, $n = 42$.
\item Energy follows power-law speed dependence: $E \propto v^{1.04}$, $r = 0.852$, $p = 1.2 \times 10^{-9}$, $n = 31$.
\item Spectral slope is speed-independent: $r = 0.193$, $p = 0.30$ (confirming $\beta_a$ reflects terrain geometry, not speed).
\item Measured $\beta_a = 1.17 \pm 0.21$, consistent with framework predictions for natural terrain with $D \approx 2.2$--$2.5$.
\end{itemize}

\noindent The TartanDrive effect sizes ($r = 0.512$--$0.880$) are substantially stronger than the LiRA-CD paved-road results ($r = 0.145$--$0.194$), confirming that the framework is most predictive for natural terrain---precisely the domain for which it is intended. This represents the first real off-road vehicle validation of the terrain-vibration spectral framework. A new subsection ``Real off-road vehicle validation: TartanDrive dataset'' has been added to the Results section.

No new simulations are required to address this comment. The TartanDrive analysis uses real measured vehicle data (not simulation). The response involves honest reframing of existing evidence, identification of validation gaps, and new real-data analysis.

\revloc{Abstract (final sentences, TartanDrive results added); Results, new subsection ``Real off-road vehicle validation: TartanDrive dataset''; Discussion (validation limitations paragraph, revised); Methods, new subsection ``TartanDrive off-road vehicle data''; Data Sources table (TartanDrive row added); Supplementary Note~13 (full TartanDrive analysis with methods, results tables, comparison, and figure).}

\bigskip

% --------------------------------------------------------------
\subsection*{Minor Comments}
% --------------------------------------------------------------

\revquote{Several abbreviations are used before being defined, or are not defined at all, e.g.\ PSD, RMS, DEM, HMMWV. Several variables are introduced without sufficient definition at first use, including $K_\sigma$ or $\omega$. This makes the mathematical development difficult to follow.}

\noindent\textbf{Response.} We thank the reviewer for identifying this significant readability issue. We have performed a comprehensive audit of every symbol, constant, variable, and abbreviation used in the manuscript and supplementary information.

The revised manuscript now includes a complete ``Notation and definitions'' paragraph at the beginning of the Results section that defines \emph{all} symbols used subsequently, organized into categories: terrain parameters, vehicle parameters, dynamic response quantities, fatigue parameters, scaling and statistics, ISO~8608 parameters (for comparison), terrain generation parameters, and numerical method quantities. Every abbreviation is expanded at first use and collected in the notation paragraph.

Specific additions and clarifications include:

\begin{itemize}[leftmargin=*]
\item All abbreviations defined: PSD, RMS, DEM, DFA, fBm, IRI, HMMWV, AIC, BIC, 2-DOF, FFT, DBC, RK4, LiDAR, 3DEP, ROC, AUC, LHS, OpenVD, WCS, DHM, GPS, CI, SD, IQR
\item $K_\sigma$: defined with units (Pa$\cdot$s$^2$/m), physical meaning (stress transfer parameter), and typical range
\item $\omega = vk$: explicitly defined as temporal angular frequency with derivation from spatial wavenumber
\item $f = \omega/(2\pi)$: temporal frequency in Hz, distinguished from angular frequency
\item $C_\omega = C_z v^{\beta_t - 1}$: temporal PSD amplitude defined and related to spatial amplitude
\item $r$ (Pearson correlation) vs $r_{\text{DS}}$ (Diamond-Square variance reduction factor): disambiguated
\item $n$ (sample size) vs $n$ (ISO spatial frequency) vs $n_i$ (cycle count): disambiguated with context notes
\item $w$ (ISO waviness exponent) vs $w(n)$ (Hann window): disambiguated with explicit note
\item $\gamma$, $\gamma_E$, $\gamma_f$, $\gamma_A$, $\gamma_B$, $\gamma_C$: all scaling exponent subscripted variants defined
\item $V \in \{0,1,2\}$: vehicle index for variance decomposition model
\item $p$: $p$-value explicitly defined (probability under the null hypothesis)
\item $F_s(t)$, $F_t(t)$: suspension force and tire force defined with equations
\item $A_{\text{eff}}$: effective cross-sectional area defined with units
\item $\sigma(t)$, $\sigma_{m,i}$: stress time history and mean stress of $i$-th cycle defined
\item $L_{\text{mission}}$: terrain traverse length defined
\item $\ddot{z}_s(t)$: sprung-mass acceleration time history defined
\item $a_z(t)$: vertical acceleration (for LiRA-CD sensor data) defined
\item $E_{\mathrm{vib}}$: total vibration energy from PSD integral defined
\item $\hat{\beta}$: estimated (fitted) spectral exponent defined
\item $\varepsilon$, $N(\varepsilon)$: DBC box size and box count defined
\item $\mathcal{N}(0,\sigma^2)$: Gaussian distribution notation defined
\item $\mathbf{x}$, $\mathbf{k}_1$--$\mathbf{k}_4$: state vector and RK4 derivatives defined
\item $\mathcal{D}_{\text{fatigue}}$: cumulative damage defined with Miner's rule formula
\item $Q = 1/(2\zeta)$: quality factor defined
\item $\Delta t$, $\Delta f$: time step and frequency resolution defined with values
\item Supplementary information includes a note directing readers to the main manuscript notation paragraph, with key definitions repeated where needed for readability. A separate ``Additional notation'' paragraph defines supplementary-only symbols ($\alpha$, $\sigma_f'$, $b$, $K_0$, $\Delta\omega_{\text{eff}}$, $H(X)$, $z_{\max}$, $z_{\min}$, $n_{\text{box}}$).
\end{itemize}

\noindent We have also addressed the $\beta$ notation transitions ($\beta_t$ for terrain PSD slope, $\beta_a$ for vehicle acceleration PSD slope, $\beta$ without subscript used only in equations where context is unambiguous) by ensuring consistent subscript usage throughout.

\revloc{Main manuscript: Results, ``Notation and definitions'' paragraph (comprehensive, covering all symbols used in the manuscript across 10 thematic categories: terrain parameters, vehicle parameters, dynamic response, fatigue parameters, scaling and statistics, ISO~8608 parameters, terrain generation, numerical methods, fractal dimension estimation, and abbreviations). Supplementary: Note~1, cross-reference statement plus new ``Additional notation'' paragraph defining supplementary-specific symbols. Throughout both documents: consistent subscript usage for $\beta_t$ vs $\beta_a$.}

\bigskip

\revquote{The notation is inconsistent throughout the formulation. In particular, the spectral exponent is initially introduced as $\beta$ and later redefined as $\beta_t$ to distinguish it from $\beta_a$, without a clear explanation of this transition.}

\noindent\textbf{Response.} We agree this was confusing. The spectral exponent $\beta$ serves double duty: it appears first in the theoretical derivation as the terrain PSD slope ($\beta_t = 7-2D$), and later as the measured vehicle acceleration PSD slope ($\beta_a$), which differs from $\beta_t$ by a factor of $\sim$0.80 due to suspension filtering (empirically: the theoretical terrain slope is $-2.0$ from $\beta_t = 7-2D$, while the measured vehicle acceleration slope is $-1.59$, giving the ratio $|-1.59|/|-2.0| = 0.795 \approx 0.80$; this attenuation occurs because the suspension transfer function acts as a low-pass filter that preferentially suppresses high-frequency spectral content, flattening the measured slope relative to the terrain input). The transition was not signposted.

We have added an explicit ``Notation convention for spectral exponents'' paragraph at the beginning of the Validation section (immediately after the Theoretical Framework), which states:

\begin{itemize}[leftmargin=*]
\item $\beta_t$: terrain elevation PSD exponent (theoretical: $\beta_t = 7-2D$)
\item $\beta_a$: vehicle acceleration PSD exponent (empirical, measured from $S_a(\omega)$)
\item $\beta$ without subscript: used only in equations where context is unambiguous (e.g., the geometric relationship $\beta = 7-2D$)
\item The relationship between the two: $\beta_a \approx 0.80 \, \beta_t$ due to suspension transfer function filtering
\end{itemize}

\noindent This clarifies why the measured slope ($-1.59$) differs from the theoretical prediction ($-2.0$) and makes the subscript convention explicit before any empirical results are presented. The comprehensive notation paragraph in the Results section also defines both $\beta_t$ and $\beta_a$ with their distinct meanings.

\revloc{Results, new ``Notation convention for spectral exponents'' paragraph (before Validation subsection, highlighted in blue); Notation paragraph (defines both $\beta_t$ and $\beta_a$).}

\bigskip

\revquote{The Introduction does not clearly outline the structure of the manuscript. As a reader, it is difficult to anticipate the logical sequence followed in the article.}

\noindent\textbf{Response.} We have added a ``Manuscript structure'' paragraph at the end of the Introduction that outlines the logical sequence:

\begin{enumerate}[leftmargin=*]
\item Results: notation $\rightarrow$ theoretical framework $\rightarrow$ Monte Carlo validation of $D$--$\beta$ $\rightarrow$ energy scaling law $\rightarrow$ fatigue scaling chain $\rightarrow$ multi-vehicle validation $\rightarrow$ real-terrain LiDAR validation $\rightarrow$ real-vehicle LiRA-CD validation
\item Discussion: ISO~8608 comparison, domain boundaries, limitations, future work
\item Methods: terrain generation, vehicle dynamics, spectral analysis, fatigue analysis, statistics
\end{enumerate}

\noindent This allows the reader to anticipate the logical progression from theory through simulation to real-data validation before encountering the technical content.

\revloc{Introduction, final paragraph ``Manuscript structure'' (highlighted in blue).}

\bigskip

\revquote{Some figure axes lack units, particularly those involving PSD and energy.}

\noindent\textbf{Response.} We have addressed this in two ways:

\begin{enumerate}[leftmargin=*]
\item A ``Units note'' has been added at the top of the Figure Legends section specifying all axis units: spatial PSD in m$^2$/(rad/m), acceleration PSD in m$^2$/s$^4$/(rad/s), vibration energy $E$ in m$^2$/s$^4$, fractal dimension $D$ and spectral exponents dimensionless.

\item The figure generation scripts (\texttt{generate\_figure1.py}, \texttt{generate\_figure3.py}) have been updated to include explicit axis labels with SI units on all panels. Revised figure files will be included with the final submission.
\end{enumerate}

\noindent We note that on log-log PSD plots, the convention of labeling axes as ``PSD [m$^2$/(rad/m)]'' and ``Wavenumber $k$ [rad/m]'' is followed, consistent with standard practice in vehicle dynamics literature (ISO~8608:2016; Bendat \& Piersol, 2011).

\revloc{Figure Legends, new ``Units note'' paragraph (highlighted in blue); figure generation scripts updated.}

\bigskip

\revquote{Figure 1 lacks clear panel labels (subfigures) in the figure itself, despite being referenced in the text.}

\noindent\textbf{Response.} Panel labels (a)--(d) have been added to the Figure~1 image file to match the caption references. This was also noted by Reviewer~1 (Formatting Issue~2) and is now corrected.

\revloc{Figure file: \texttt{figures/Figure1.png} updated with panel labels.}

\bigskip

\revquote{The selected values $D=2.1, 2.3, 2.5$ in Figure 1c--d are not justified. The authors should explain why these values were selected and how they relate to the full range used in the simulations and plotted in Figure 1a--b ($D=2.0$--$2.6$).}

\noindent\textbf{Response.} These three values were selected to span the simulation range ($D = 2.05$--$2.45$) at equal intervals, representing smooth ($D = 2.1$), moderate ($D = 2.3$), and rough ($D = 2.5$) terrain classes. Three curves suffice for visual clarity while demonstrating the monotonic change in spectral slope across the tested range. We have added a clarifying sentence to the Figure~1 legend: ``These values span the simulation range ($D = 2.05$--$2.45$) at equal intervals, representing smooth, moderate, and rough terrain classes.''

\revloc{Figure Legends, Figure~1 panel~(c) description (highlighted in blue).}

\bigskip

\revquote{Vehicle definitions A, B, and C are introduced in the early Results section before their physical parameters are specified. In addition, it is unclear whether the same vehicle definitions are used consistently throughout the manuscript, as different parameter sets appear to be associated with these labels in later sections.}

\noindent\textbf{Response.} This ambiguity arose because the manuscript uses two vehicle configurations: (1)~the HMMWV quarter-car ($m_s = 850.5$~kg, $k_s = 36{,}300$~N/m, $f_n = 1.04$~Hz) for the 500-simulation amplitude-coupling analysis (Figure~4), and (2)~Vehicles A, B, C (800/1000/1200~kg) for the 1500-simulation multi-vehicle validation (Figures~2--3). We have made two clarifying additions:

\begin{itemize}[leftmargin=*]
\item At first use of Vehicles A/B/C in the energy scaling section (Figure~3 description): added forward reference ``full parameter definitions in `Terrain-dominated scaling' below and Methods.''
\item At the full parameter definition: added explicit statement that ``These definitions are used consistently throughout the manuscript; all references to Vehicles A, B, and C use these parameter sets. The HMMWV configuration is used only for the 500-simulation amplitude-coupling analysis (Figure~4) and is identified explicitly.''
\end{itemize}

\revloc{Results, Figure~3 paragraph (forward reference); Results, ``Terrain-dominated scaling'' paragraph (consistency statement). Both highlighted in blue.}

\bigskip

\revquote{Latin Hypercube sampling is mentioned in the extended vehicle ensemble section but is not briefly explained.}

\noindent\textbf{Response.} We have added a parenthetical definition: ``Latin Hypercube sampling (LHS), a space-filling experimental design that divides each parameter range into $N$ equal strata and samples exactly once from each stratum, ensuring uniform coverage of the multi-dimensional parameter space without clustering.'' A reference to McKay et al.\ (1979) has been added.

\revloc{Results, ``Extended vehicle ensemble'' paragraph (highlighted in blue); Bibliography: added McKay1979.}

\bigskip

\revquote{Several scatter plots, particularly Figure 5c and Figure 6b, show substantial dispersion. The authors should justify the use of linear regression and discuss whether alternative fitting strategies or uncertainty quantification approaches were considered.}

\noindent\textbf{Response.} We have added a ``Regression methodology and dispersion'' paragraph addressing this systematically:

\begin{itemize}[leftmargin=*]
\item \textbf{Figure~5b ($\gamma$ vs $f_n$):} Linear regression is justified because the theoretical spectral integration predicts monotonic, approximately linear frequency dependence over the tested range. Residual analysis confirms homoscedasticity (residual SD = 0.066).
\item \textbf{Figure~5c ($\gamma$ vs mass):} The dispersion \emph{is} the result---the near-zero $R^2 = 0.000$ demonstrates that mass does not predict the exponent. Alternative models (quadratic, exponential) yield no improvement.
\item \textbf{Figure~6b (Copenhagen):} The dispersion reflects the physical spatial-scale mismatch (micro-texture dominating vehicle excitation on paved roads), not a statistical modeling deficiency. The near-zero/wrong-sign correlation is the scientifically meaningful finding (domain boundary).
\end{itemize}

\revloc{Results, new ``Regression methodology and dispersion'' paragraph after the extended ensemble section (highlighted in blue).}

\bigskip

\revquote{Figure 5f is not discussed in the text.}

\noindent\textbf{Response.} We have added a sentence in the ``Regression methodology and dispersion'' paragraph: ``Figure~5f visualizes the joint distribution of all 100 vehicles in the 3D parameter space (mass, frequency, damping), confirming that terrain geometry drives scaling consistently across the entire design space without parameter-dependent clustering.''

\revloc{Results, ``Regression methodology and dispersion'' paragraph (final sentence, highlighted in blue).}

\bigskip

\revquote{Some results are repeated several times, e.g., $E \propto (D-2)^{-3.1}$. Once introduced, this relation should be referenced rather than repeatedly restated.}

\noindent\textbf{Response.} We have revised the Discussion to eliminate redundant restatement of the energy scaling relationship. The first occurrence in Results (Figure~3 description, with full statistics) remains as the definitive statement. Subsequent mentions in the Discussion now use back-references (``the energy--$D$ correlation observed in Figure~3'') rather than restating the equation. The two-parameter model $E \propto C_z^{0.94} \times \beta_a^{-0.09}$ is similarly stated once in full (Results) and referenced thereafter.

\revloc{Discussion, first two paragraphs condensed and cross-referenced (highlighted in blue).}

\bigskip

\revquote{The manuscript contains formatting artifacts, including markdown-style bold markers such as \texttt{**Critical methodological note:**}.}

\noindent\textbf{Response.} We have searched the entire manuscript and supplementary for markdown-style formatting artifacts (\texttt{**text**}, \texttt{\_text\_}, etc.). No instances were found in the current revision. All bold text uses \texttt{\textbackslash textbf\{\}} and all italics use \texttt{\textbackslash textit\{\}} as required by \LaTeX. If the reviewer observed these artifacts in the original submission, they have been corrected during the revision process.

\revloc{Throughout manuscript (verified, no markdown artifacts present).}

\bigskip

\revquote{Some references appear to be missing or incomplete (e.g., see line 331).}

\noindent\textbf{Response.} The erroneous ``[?]'' citation symbols at line~331 (and line~731) arose from two references (Cover \& Thomas 2006; Musgrave et al.\ 1989) that were cited in the text but missing from the bibliography. Both entries have been added and now resolve correctly. This was also identified by Reviewer~1 (Formatting Issue~1). We have verified that all \texttt{\textbackslash cite\{\}} commands resolve to valid bibliography entries in the revised manuscript.

\revloc{Bibliography: added \texttt{Cover2006} and \texttt{Musgrave2000} entries; added \texttt{McKay1979} for Latin Hypercube.}


\bigskip\noindent\rule{\textwidth}{0.4pt}

% ==============================================================
\section*{Reviewer \#3}
% ==============================================================

\noindent Reviewer \#3 raised four major issues. We address each below.

\bigskip

% --------------------------------------------------------------
\subsection*{Major Issue 1 --- Amplitude-complexity artifact}
% --------------------------------------------------------------

\revquote{The negative correlation between vibration energy and fractal dimension ($E \propto (D-2)^{-3.1}$) appears to be an artifact of the Diamond-Square terrain generation algorithm rather than a universal physical law. This undermines the central claim that fractal dimension predicts fatigue.}

\noindent\textbf{Response.} The reviewer is correct, and we have made this point explicit throughout the revised manuscript. The negative $E$--$D$ correlation is indeed an artifact of the Diamond-Square algorithm's amplitude-complexity coupling ($C_z \propto (D-2)^{-3.12}$, $R^2 = 0.94$), not a universal physical law. We now state this unambiguously in multiple locations: the Results text, Figure~3 caption, Discussion, and a dedicated ``Critical methodological note'' paragraph.

However, the artifact does not undermine the scientific contribution. The key findings are:

\begin{enumerate}[leftmargin=*]
\item The geometric relationship $\beta_t = 7-2D$ (derived from self-affine theory, not fitted to data) correctly predicts spectral slope from fractal dimension. This is validated on synthetic terrain ($r = -0.956$, $n = 1500$) and real USGS LiDAR terrain ($r = -0.62$, $p < 0.05$).

\item The two-parameter framework $(C_z, \beta)$ explains 96\% of energy variance, demonstrating that \emph{both} amplitude and spectral structure are needed for complete terrain characterization.

\item The constant-amplitude decoupling experiment ($R^2 = 0.834$, $n = 449$) demonstrates that spectral structure independently affects vibration energy when amplitude is controlled. Energy \emph{increases} with $D$ at constant amplitude (positive slope $= +2.96$), opposite to the artifact-driven coupled case. This proves that $\beta$ carries independent physical information beyond what amplitude alone provides.

\item The fatigue scaling cascade (2:1 ratio between $\gamma_f$ and $\gamma_E$ across 100 vehicles) validates the multi-physics chain regardless of the sign of the $E$--$D$ relationship.
\end{enumerate}

\noindent The revised manuscript clearly separates the artifact (negative $E$--$D$ correlation due to amplitude coupling in synthetic terrain) from the science (two-parameter framework, geometric $\beta = 7-2D$ relationship, independent spectral structure effect, fatigue scaling cascade). Natural terrain may exhibit different or opposite amplitude-complexity relationships; the framework applies regardless.

\revloc{Results ``Energy scaling law,'' Critical methodological note paragraph; Figure~3 caption (explicit artifact statement); Discussion; Supplementary Note~9 (full analysis of coupling).}

\bigskip

% --------------------------------------------------------------
\subsection*{Major Issue 2 --- Real terrain validation weakness}
% --------------------------------------------------------------

\revquote{Validation using real data is limited: only 13 USGS tiles with moderate correlation ($r = -0.62$), and the LiRA-CD analysis shows only weak effect sizes ($r \approx 0.15$--$0.19$). This is insufficient to support the broad claims made.}

\noindent\textbf{Response.} We agree that the real-data validation is preliminary rather than comprehensive, and have revised the manuscript to frame it accordingly. Several points of clarification:

\textbf{USGS 3DEP (13 tiles, $r = -0.62$, $p < 0.05$):} This validates the $D \rightarrow \beta_t$ link on real terrain. The moderate correlation compared to simulation ($r = -0.956$) is expected: real terrain combines multiple geological processes, exhibits multifractal properties, and departs from ideal fBm. The correlation is statistically significant despite small $n$, confirming that the theoretically predicted negative relationship exists in independently measured natural terrain. Expansion to 25 class-balanced tiles is in progress to improve statistical robustness.

\textbf{LiRA-CD (8,609 segments, $r = 0.15$--$0.19$, $p < 10^{-40}$):} This validates the $\beta_a \rightarrow E$ link on real vehicle data with a large sample size. The modest effect size reflects real-world complexity (non-terrain vibration sources, GPS matching noise, speed variation, consumer vs military vehicle), not framework failure. The extremely low $p$-values confirm these are genuine physical relationships. Importantly, the LiRA-CD validates the spectral relationship \emph{between} terrain roughness (IRI) and vehicle vibration spectra/energy---precisely the link that the framework predicts---using 8,609 independent measurements.

\textbf{Together:} The USGS analysis validates $D \rightarrow \beta_t$ on real terrain geometry; LiRA-CD validates $\beta \rightarrow E$ on real vehicle dynamics. These are independent datasets using independent methods in independent geographic settings, making their combined confirmation robust.

\textbf{NEW: TartanDrive off-road validation.} We have now completed analysis of the TartanDrive dataset (Triest et al., ICRA 2022): a Yamaha Viking ATV traversing diverse natural terrain with IMU at 200~Hz. Across 42 trajectory segments (9 traversals, 10-second windows), we find:

\begin{itemize}[leftmargin=*]
\item $E$ vs $\beta_a$: $r = 0.512$, $p = 5.4 \times 10^{-4}$ --- spectral slope predicts vibration energy on real off-road terrain.
\item $E \propto v^{1.04}$: $r = 0.852$, $p = 1.2 \times 10^{-9}$ --- power-law speed dependence confirmed (theory: $E \propto v^{6-2D}$; observed slope implies $D_{\mathrm{eff}} \approx 2.48$).
\item $\beta_a$ is speed-independent ($r = 0.193$, $p = 0.30$), confirming spectral slope reflects terrain geometry.
\item Measured $\beta_a = 1.17 \pm 0.21$ falls within framework predictions for natural terrain.
\end{itemize}

\noindent Critically, the TartanDrive off-road effect sizes ($r = 0.512$--$0.880$) are \emph{substantially} stronger than LiRA-CD paved-road results ($r = 0.145$--$0.194$). This is consistent with the framework's design intent and directly addresses the reviewer's concern: the framework is most predictive precisely for the natural-terrain domain where it is intended to apply.

\textbf{Revised framing:} All claims are now moderated to distinguish ``demonstrates the existence and direction of predicted relationships'' from ``proves quantitative prediction accuracy.'' The framework is presented as establishing scaling relationships and providing geometric interpretation, not as an operational prediction tool ready for deployment.

\revloc{Abstract; Discussion (validation paragraph); Real terrain comparison section; LiRA-CD section.}

\bigskip

% --------------------------------------------------------------
\subsection*{Major Issue 3 --- Universality contradiction}
% --------------------------------------------------------------

\revquote{The title claims universality but the framework explicitly fails on paved roads. The term ``universal'' is contradicted by the Copenhagen negative result.}

\noindent\textbf{Response.} We fully agree. The title has been changed from ``Universal spectral scaling of terrain-induced fatigue'' to ``Spectral scaling of natural-terrain-induced fatigue.'' The word ``universal'' has been removed from the entire manuscript. The abstract now ends with an explicit domain boundary statement: ``This framework enables quantitative prediction of vibration severity and fatigue risk for natural terrain where digital elevation models are available; it is not applicable to engineered paved surfaces.''

The Copenhagen negative result is retained as a scientifically valuable boundary condition that defines where the methodology applies (natural terrain with macro-topographic relief) and where it does not (paved roads where centimeter-scale pavement micro-texture dominates). This is analogous to other physical theories with well-defined applicability limits (e.g., Newtonian mechanics at $v \ll c$, ideal gas law at low pressures).

No new simulations are required. This is a framing correction.

\revloc{Title; Abstract; Discussion (domain boundary paragraph retained, universality claims removed).}

\bigskip

% --------------------------------------------------------------
\subsection*{Major Issue 4 --- Spatial scale limitation}
% --------------------------------------------------------------

\revquote{The methodology is limited to terrain features at DEM-measurable scales (meters to kilometers), while real vehicle vibration on roads involves centimeter-scale features. This spatial scale limitation is not adequately discussed.}

\noindent\textbf{Response.} We agree this limitation needed more thorough discussion and have added Supplementary Note~8A (``Domain of Applicability and Spatial Scale Limitations'') providing a comprehensive analysis of the two terrain excitation regimes:

\textbf{Regime 1 (Natural terrain):} Vehicle excitation dominated by macro-topographic geometry at 1--1000~m wavelengths (hillslopes, rock outcrops, erosion channels). These features are captured by DEMs at 1--30~m resolution. The methodology works in this regime ($r = -0.583$, $p = 0.036$ on USGS terrain).

\textbf{Regime 2 (Paved roads):} Vehicle excitation dominated by pavement micro-texture at 0.001--0.5~m wavelengths (asphalt roughness, potholes, cracks, joints). These features are not captured by any available DEM. The methodology fails in this regime ($r \approx 0$, wrong direction, Copenhagen).

The spatial scale mismatch is now explicitly framed as an operating boundary rather than a limitation that could be overcome with higher-resolution DEMs. The key evidence: increasing DEM resolution from 30~m to 1~m does \emph{not} improve the correlation (it weakens from $r = +0.102$ to $r = +0.037$), proving that the relevant roughness features are at sub-DEM scales.

\textbf{Practical implication:} The methodology applies to its intended domain (military off-road operations, agricultural terrain, autonomous off-road navigation, planetary exploration) but not to highway maintenance or urban route optimization. This is now stated explicitly in both the main manuscript and supplementary.

\revloc{Main manuscript: Figure~6 and surrounding text; Discussion (spatial scale paragraph). Supplementary: new Note~8A ``Domain of Applicability and Spatial Scale Limitations'' (comprehensive analysis with table).}


\bigskip\noindent\rule{\textwidth}{0.4pt}

\section*{Summary of Changes}

\begin{enumerate}[leftmargin=*]
\item Title changed to ``Spectral scaling of natural-terrain-induced fatigue''
\item Abstract: scoped to natural terrain, domain boundary explicit, variance values data-verified, TartanDrive results added
\item Introduction: ISO~8608 comparison paragraph added; lifecycle cost claim softened; manuscript structure roadmap added
\item Results: comprehensive notation paragraph; $\beta_t$/$\beta_a$ notation convention paragraph; bridging paragraph for fatigue section; critical methodological notes on amplitude-complexity artifact; \textbf{new subsection ``Real off-road vehicle validation: TartanDrive dataset'' with $n=42$ real off-road segments}
\item Discussion: ISO~8608 comparison table; speed dependence paragraph; role of $\beta$ paragraph with constant-amplitude results; validation limitations honestly stated
\item Methods: ``Scope of fatigue analysis'' paragraph added
\item Throughout: ``vehicle-independent'' $\rightarrow$ ``terrain-dominated''; ``fatigue prediction'' $\rightarrow$ ``fatigue-risk scaling''; ``universal'' removed; $r$ $\rightarrow$ $r_{\text{DS}}$ for Diamond-Square variance reduction factor
\item Supplementary: new Note~8A (spatial scale limitations), new Note~12 (constant-amplitude decoupling experiment), cross-reference to main notation paragraph
\item Figures: Figure~1 panel labels added; Figure~3 corrected (energy-vs-$D$ plot, slope $-3.11$); axis units explicitly labelled in all figure scripts; Units note added to Figure Legends; \textbf{new TartanDrive validation figure (Supplementary Figure~6)}
\item Bibliography: missing entries added (Cover~2006, Musgrave~1989, Triest~2022)
\item New simulation: constant-amplitude decoupling ($R^2 = 0.834$, 449 simulations) proving independent $\beta$ effect
\item \textbf{New real-data validation: TartanDrive off-road ATV dataset ($E$ vs $\beta_a$: $r=0.512$, $p<10^{-3}$; $E \propto v^{1.04}$: $r=0.852$, $p<10^{-9}$)}
\end{enumerate}

\end{document}
